\chapter{Introduction}

\section{Introduction}
Complex exponentials are functions of particular interest to mathematicians, scientists, and engineers. Recall the following:

\begin{equation}
	\frac{d}{dt} e^{st} = s e^{st}
\end{equation}

\begin{equation}
	\int e^{st} dt = \frac{1}{s} e^{st} + C
\end{equation}

where $s$ is a complex constant.

The complex exponential $e^{st}$ is an eigenfunction of differentiation with eigenvalue $s$. Because of this fact, complex exponentials are solutions to many differential equations and generate simple input-output relationships for systems described by differential equations.

The preceding commentary might suggest that systems described by differential equations could be simply analyzed by decomposing inputs into a sum of complex exponentials. This possibility was first explored by Joseph Fourier in 1821 as he used complex exponentials to solve problems involving heat flow. This type of analysis, called Fourier analysis in his honor, is used today in diverse fields to great effect.
In this chapter, we explore the two primary ways to decompose functions into sums of complex exponentials. These are the Fourier series and transforms. In Chapter 3 we will explore which types of systems can be usefully analyzed with Fourier analysis.
